
\section{The Congregations’ Right to Call and the Refusals of Ordination}

“Folkebladet,” February 6, 1895.

To deprive the congregations of their right to call is so old an offense within the Church that it is no wonder if it returns again whenever there is occasion.

For there is a tendency, which is in the blood of us all, that the freedom of others is hard to endure, if in any way it restricts our own and runs counter to our personal advantage.

If a tightly united, carnal clergy could get matters arranged in such a way that congregation was merged into pastorate and these pastorates were “good” or “fat,” and then the pastors could divide them among themselves, without these wronged congregations having anything to do with the matter, is there any doubt that all carnal pastors would find such an arrangement exceedingly good?

And since there are always many purely carnal pastors, and since all pastors without exception are not exclusively Spirit and therefore have their temptations to lord it over the Lord’s heritage, it is continually necessary that the congregations be on their guard, lest they lose the beautiful jewel of liberty and become a livelihood for those who desire to divide the offices among themselves.

Where the pastors have become generally spiritually lifeless in their understanding of the congregation and of their position in it, this will make itself known in various ways. Secret pastoral conferences will be formed and maintained, where the class interests of the clergy may be advanced; secret pastoral conferences have never been anything else, nor can they become it.

Next, the various attempts will be made to dictate to the congregations which pastor they are to have and not have. If it is a mission congregation, it is advised to surrender the right to call to the mission committee; if it is a more independent congregation, then it is to “apply to the chairman” and get his advice. If the congregation in particular wishes to have a man with whom the secret council is not pleased, then it is said: “No, he is not to be had; there is no use in calling him.” In short, there are many crooked ways by which the congregations’ free right to call is placed in danger once a carnal mind and machine politics first creep into the Church.

In the United Church we see things clearly enough. On the one side there are the “appointments”; on the other side there are the “refusals of ordination.” Both are the fruit of the same bitter root.

There is more than one mission congregation which has been advised to surrender its right to call to the mission committee, and there is more than one congregation which has followed the advice and has been “appointed over.” In the “Luth. Kirkeblad” we even find the following from a congregation in the far West:

“Likewise we thank the mission committee for its labor with the appointment and oversight of mission pastors, and with prayer that they will as soon as possible send us a pastor again.” (“L. K.” 1895, p. 26.)

Yes, that is the way it goes. State-church order and state-church rule are smuggled in upon the congregations straight against God’s Word and the freedom and right of the congregations.

Likewise with the refusals of ordination; first people are made to imagine that ordination is a priestly matter or an episcopal act, and then ordination is denied on purely church-political grounds, and then the congregations are to discover that they have come into a predicament from which they cannot get out.

It is first an affront to the congregations’ understanding of God’s Word and Christian order. For where the congregations are the least little bit instructed, they no longer have the Catholic superstition that ordination is an act of the assembled clergy, but rather a Christian act which can be carried out by any Christian congregation, just as the whole Lutheran Church unanimously confesses in the Smalcald Articles, where we read the following words, which may also have their application to our own circumstances with suitable changes.

“But since by divine right there is now no difference between bishops, pastors, or priests, it is beyond all doubt that when a pastor in his congregation ordains capable persons to church offices, then this ordination is valid according to divine right. Therefore, when the bishops persecute the Gospel and refuse to ordain capable men, then in such a case every congregation has full right and warrant to ordain its own church servants. For where the Church is, there is also the command to preach the Gospel. Therefore the congregations must retain the power to call, choose, and ordain church servants.”

So broad-minded one once was in the Lutheran Church; but now men will have the right to “refuse to ordain capable men,” and those who ordain them they will make into rebels and transgressors of the law. First they will make people superstitious, and then use the superstition as a chain upon the congregations, so that they shall not dare to exercise their right to call for fear of not getting their called pastor ordained.

Such abuse goes on for a time, especially among dead and sleeping people. We know, however, that it is dangerous work to build church bodies upon ignorance, superstition, and bondage. The United Church will find that it harms itself by not taking a clear and straightforward stand upon the ground of free-church congregational liberty and fully respecting the congregations’ right to call.

Georg Sverdrup